\section{Discussion}
\label{sec:discussion}

Symmetric bodies are ahead from the first generation of controller search, under every controller that has a shared clock, and learning does not create the lead; it shrinks it in relative terms without closing it (\cref{sec:r2}). Two controls place most of the lead in the bodies. Tree bodies chosen to have the resting shape of symmetric bodies close about half of it, and bodies that duplicate modules without mirroring them close most of it with a different, tree-like shape (\cref{sec:r3}). In evolution, tied bodies match symmetric ones as well (\cref{tab:e1}). What these bodies share is not the mirror plane but what a generative constraint does to the distribution of bodies: duplicated limbs, more of which touch the floor, and, for mirroring specifically, low and sprawled bodies whose legs sweep horizontally.

\caveat{``Most'' and ``about half'' should follow the final CIs (\cref{tab:controls}). If v2 shrinks the tied effect, fall back to: symmetric bodies start ahead, partly through shape, with duplication as a suggestive result.}

Symmetric bodies do move differently: they alternate their mirrored limbs and lock them in phase, which tree bodies almost never do (\cref{sec:r4}). It would be natural to conclude that symmetry helps because it supports a coordinated gait that a controller can find easily. Three observations argue against this as the main mechanism. Tied bodies reach symmetric scores without alternating or locking. The trained gaits of locked symmetric bodies depend on their timing as much as any gait does, but they are no more sensitive to the timing between their two sides than to any other timing of the same number of hinges, whether the lag is shifted or the lock is broken by detuning (\cref{sec:r5}). And the advantage is present before learning, when no gait has yet been found. The lock tends to accompany good gaits on symmetric bodies, including across training seeds of the same body, but these gaits do not depend on it.

\caveat{Scope: the phase-shift test addresses whether trained gaits depend on the lock, not whether the lock helps the search find them. Keep ``mechanism'' tied to the trained gait, or run the constrained-training test suggested in \cref{sec:r5}.}

A shared clock is necessary: with a frequency per hinge every body collapses, and the symmetric advantage disappears with the scores (\cref{sec:r6}). Given a shared clock, the advantage appears both with an open-loop controller with explicit per-hinge phases and with a decentralised neural controller whose parameter count does not depend on the body. In evolution the neural controller shows the advantage only with a larger inner budget (\cref{sec:r1}), although on fixed bodies it is present at the small budget too.

Comparisons of morphological encodings usually credit the regularity an encoding imposes (symmetry, repetition, modularity) for the performance of the bodies it produces. Our results suggest that a large part of such an advantage can come from the body statistics the encoding changes: how many modules are duplicated, how many limbs reach the floor, and the resting shape. Controls that match these statistics, here a tied genome and shape-matched bodies, are cheap and should be standard before an advantage is attributed to the regularity itself.

The study has several limitations. The central symmetric $-$ tree gap on fixed bodies was only marginally significant with 40 bodies per genome (\cref{tab:e2}), and the shares of the gap closed by the controls have wide intervals; the v2 runs double the bodies. We study one task, forward locomotion on flat ground, with one set of modules; the hinge grid tests robustness to the actuator, but per-setting gaps are noisy. The fixed-body experiments use random bodies from each genome's initial distribution, and evolved bodies may differ. Shape matching is partial: tree bodies always reach the module cap and stay taller than symmetric bodies even after matching, and although regression adjustment for the remaining mismatch leaves the residual gap unchanged, unmeasured shape differences remain. Gait classes come from an automatic detector validated on a different module set. Finally, the absence of a \dann\ gap at the small inner budget in evolution is not yet explained.
\todo{Update the limitations after v2.}

Matching tree bodies to tied bodies on leg hinges (E5) will show whether duplication works by adding usable legs. A dose--response over the strength of a symmetry or duplication constraint, and other regularities (repetition, segmentation, radial symmetry), would test whether the question ``which body statistics does an encoding change?'' predicts its advantage more generally.
